\section{Aplicaciones y resultados numéricos}
\label{bre:res:section}

Las aplicaciones distinguen la obtención de una raíz algebraica, la
construcción de una condición inicial y la clasificación de la trayectoria
resultante. El cierre armónico se evalúa mediante los residuos de la
sección~\ref{bre:bal:sec}; las integraciones y los contactos se interpretan
con los criterios de la sección~\ref{bre:num:seccion}. En las comparaciones
fraccionarias, cada reinicio tiene terminal inferior \(t_0=0\) y memoria
inicial nueva, salvo la continuación expresamente descrita con historia
heredada. Las frecuencias espectrales se expresan en ciclos por unidad de
tiempo; las frecuencias \(\omega\) del balance son angulares.
Las raíces y magnitudes calculadas se muestran redondeadas; las identidades
exactas corresponden a sus expresiones algebraicas. Los parámetros
decimales declarados fijan el problema de cada ensayo.

\subsection{Chua no suave de orden entero}
\label{bre:res:chua-integer}

Se considera el campo de Chua de \eqref{bre:bal:chua}, con
\[
 \begin{aligned}
 (\alpha,\beta,\gamma)&=(8.4562,12.0732,0.0052),\\
 (m_0,m_1)&=(-0.1768,-1.1468),\qquad q=1,
 \end{aligned}
\]
y característica \(h(x)=m_1x+(m_0-m_1)\operatorname{sat}(x)\)
\cite{LeonovKuznetsov2013}. Para \(\beta+\gamma\ne0\), las ecuaciones de
equilibrio segunda y tercera se resuelven como
\begin{equation}
 y=\frac{\gamma x}{\beta+\gamma},\qquad
 z=-\frac{\beta x}{\beta+\gamma},\qquad
 h(x)+\frac{\beta x}{\beta+\gamma}=0.
 \label{bre:res:chua-equilibrium-reduction}
\end{equation}
La última identidad se obtiene sustituyendo las primeras en
\(y-x-h(x)=0\). Sean
\(\kappa=m_1+\beta/(\beta+\gamma)\) y \(d=m_0-m_1\).
En \(|x|\leq1\), queda \((\kappa+d)x=0\); como
\(\kappa+d>0\), sólo se obtiene el origen. En \(|x|>1\),
\(\kappa x+d\operatorname{sign}(x)=0\), por lo que
\[
 x=-\frac d\kappa\operatorname{sign}(x).
\]
Aquí \(\kappa<0\) y \(-d/\kappa>1\): hay exactamente dos raíces externas,
una por signo, cuyas otras coordenadas da
\eqref{bre:res:chua-equilibrium-reduction}. Ninguna está en un quiebre.
Para una característica diferenciable en el equilibrio, con pendiente
\(m=h'(x_e)\), el Jacobiano y su polinomio son
\begin{equation}
 \begin{gathered}
 J_e=\begin{pmatrix}
 -\alpha(1+m)&\alpha&0\\1&-1&1\\0&-\beta&-\gamma
 \end{pmatrix},\\
 p_e(s)=[s+\alpha(1+m)]D(s)-\alpha(s+\gamma),\\
 D(s)=s^2+(1+\gamma)s+\beta+\gamma.
 \end{gathered}
 \label{bre:res:chua-jacobian}
\end{equation}
La expansión del determinante sigue \eqref{bre:bal:characteristic}; se usa
\(m_0\) en el origen y \(m_1\) en los equilibrios externos. La misma
matriz permite aplicar Matignon al cambiar el orden, sin cambiar los
equilibrios del campo.

La raíz de fase, la ganancia y la amplitud
calculadas por \eqref{bre:bal:phase-polynomial} y
\eqref{bre:bal:saturation-df} son
\[
 \begin{aligned}
 \omega_0&=2.039186939959,\\
 k&=0.209867354515,\qquad A_0=5.856145086252.
 \end{aligned}
\]
La reconstrucción modal de \eqref{bre:bal:seed} produce
\begin{equation}
 X_0=\begin{pmatrix}
 5.856145086252\\0.369331578246\\-8.366536168324
 \end{pmatrix}.
 \label{bre:res:chua-integer-seed}
\end{equation}
Los residuos del cierre complejo y de la función descriptiva fueron,
respectivamente, \(3.26\times10^{-15}\) y \(1.98\times10^{-13}\). Estas
cifras verifican las ecuaciones usadas para generar la semilla; la
continuación y los diagnósticos posteriores determinan su destino.
Una búsqueda independiente por recurrencias, con 512 condiciones iniciales,
identificó dos conjuntos recurrentes y 22 contactos. La fracción
\(22/512=0.04296875\) se refiere al dominio y a la muestra elegidos.

La comparación de vecindades empleó las mismas 504 coordenadas en dos
integraciones independientes: tres equilibrios, siete radios entre
\(10^{-5}\) y \(10^{-2}\), y 24 puntos uniformes en cada bola. El horizonte
fue \(T=500\), con descarte inicial de 120; las tolerancias de clasificación
fueron \(10^{-3}\) para equilibrio y \(0.08\) para la distancia normalizada
a la nube. Coincidieron las 504 clases: 152 trayectorias hacia el origen,
198 acotadas no triviales y 154 divergentes según el detector. Ninguna
alcanzó la nube candidata. Los espectros de Lyapunov finitos fueron
\[
 \begin{aligned}
 \widehat\lambda^{(1)}&=(0.14018,0.01319,-1.20304),\\
 \widehat\lambda^{(2)}&=(0.14591,-0.00052,-1.17115).
 \end{aligned}
\]
El exponente máximo positivo y la geometría de la
Fig.~\ref{bre:res:fig-chua-integer} respaldan un candidato caótico bajo esas
ventanas. La ausencia de contactos sólo establece compatibilidad con
ocultedad en el sondeo finito.

El control EFORK--3 de referencia utilizó un protocolo diferente:
\(T=320\), descarte de 80, otras 504 coordenadas y un clasificador por huellas
de sección. Produjo 369 convergencias a equilibrio, 69 divergencias y 66
trayectorias hacia otro destino, también sin contactos. Estos conteos no se
suman a los anteriores ni se interpretan como repeticiones del mismo ensayo.

\begin{figure*}[t]
 \centering
 \includegraphics[width=0.72\textwidth]{chua_entero.png}
 \caption{Trayectoria candidata del Chua no suave entero, localizada mediante
 semilla armónica y continuación. La representación corresponde a una ventana
 postransitoria; la clasificación dinámica utiliza además exponentes y sondeos.}
 \label{bre:res:fig-chua-integer}
\end{figure*}

\subsection{Chua de Caputo no suave: continuación de una semilla sesgada}
\label{bre:res:chua-pwl}

Se mantienen los cinco parámetros anteriores y se fija \(q=0.9998\).
Como control, iniciar directamente el PVI de Caputo desde
\eqref{bre:res:chua-integer-seed}, con ABM de memoria completa,
\(h=0.01\), \(T=500\) y descarte de 250, produjo una respuesta acotada
con frecuencia dominante \(0.3679852806\). Las medianas de la prueba 0--1
fueron
\[
 \begin{aligned}
 K_x&=0.0736852160,\qquad K_y=0.0214448125,\\
 K_z&=-0.0012988924.
 \end{aligned}
\]
Es una reconstrucción numérica relacionada con el caso de
Danca~\cite{Danca2017}; la condición inicial se obtuvo por balance y no se
atribuye a la fuente como dato publicado.

La localización sesgada resolvió el cierre de
\eqref{bre:bal:bias-closure}. Se ensayaron 2430 inicios y se seleccionó,
en el par nominal de pendientes,
\[
 \begin{aligned}
 A&=4.577882721023,&c&=2.776397003274,\\
 \omega&=2.040286051079,&N_1&=0.210022792962,\\
 \Psi_0&=0.408770376181.
 \end{aligned}
\]
La integración trapezoidal periódica de los coeficientes utilizó 8192 nodos
y dio un residuo de \(4.55\times10^{-16}\). Al reevaluar las integrales
continuas con cuadratura adaptativa y separar los quiebres de la saturación,
el residuo fue \(1.9\times10^{-7}\), menor que el umbral \(10^{-4}\). La
segunda cifra controla la aproximación de las integrales, no sólo la
resolución del sistema discretizado. La semilla y el extremo obtenido son
\begin{equation}
 \begin{aligned}
 X_{\rm sem}&=\begin{pmatrix}
 7.354279724297\\0.290968778784\\-9.316054818666
 \end{pmatrix},\\
 X_{\rm end}&=\begin{pmatrix}
 0.822310709850\\1.465272703995\\1.292308250248
 \end{pmatrix}.
 \end{aligned}
 \label{bre:res:pwl-seeds}
\end{equation}
La homotopía \eqref{bre:bal:homotopy} se integró en 21 niveles,
\(\varepsilon_k=0.05k\), con \(h=0.01\) y 60 unidades por nivel. Cada nivel
conservó 30 unidades después de su tramo inicial de ajuste. Las sumas de
Caputo mantuvieron la historia desde el primer nivel, como exige
\eqref{bre:num:continuacion-integral}; el tiempo total fue \(21\cdot60=1260\).
La Fig.~\ref{bre:res:fig-pwl-continuation} muestra ese transporte.

Desde \(X_{\rm end}\) se inició un PVI autónomo nuevo, con \(h=0.01\),
\(T=300\) y descarte de 100. Su frecuencia dominante fue
\(0.3699815009\), y sus entropías espectrales
\((0.11526785,0.10497088,0.10339964)\). La respuesta tiene apariencia
periódica en la ventana retenida, como muestra la
Fig.~\ref{bre:res:fig-pwl-final}; no se afirma una solución periódica exacta
del PVI de Caputo con terminal finito. Las 36 perturbaciones de
\(X_{\rm end}\), doce para cada radio \(10^{-6},10^{-4},10^{-2}\),
alcanzaron la nube objetivo con \(h=0.01\) y memoria completa. Este resultado
mide sensibilidad al dato inicial a paso fijo. El refinamiento
\(h,h/2,h/4\) de ese banco sigue pendiente.

\begin{figure*}[t]
 \centering
 \includegraphics[width=0.96\textwidth]{continuacion_caputo.pdf}
 \caption{Continuación causal del Chua no suave, \(q=0.9998\).
 Los 21 niveles de la homotopía comparten una sola historia. El estado final
 se utiliza después para definir un PVI autónomo con memoria reiniciada.}
 \label{bre:res:fig-pwl-continuation}
\end{figure*}
\begin{figure*}[t]
 \centering
 \includegraphics[width=0.96\textwidth]{chua_caputo_pwl.pdf}
 \caption{Respuesta postransitoria desde \(X_{\rm end}\),
 \(q=0.9998\), \(h=0.01\). La regularidad observada y los contactos desde
 perturbaciones son resultados finitos; no establecen periodicidad exacta.}
 \label{bre:res:fig-pwl-final}
\end{figure*}

El sondeo de vecindades utilizó bolas uniformes, con la construcción de la
sección~\ref{bre:num:sondeos}. Cada sonda se integró con \(h=0.01\),
\(T=200\), descarte de 50 y umbral de contacto \(D_{90}\leq0.5\).
El detector de equilibrio requería, desde \(t=25\), distancia menor que
\(10^{-3}\), norma del campo menor que \(10^{-4}\) y 50 pasos consecutivos;
al final del horizonte se admitía distancia a equilibrio no mayor que 0.5.
El detector de divergencia usó cinco normas consecutivas mayores que 80 o
cinco razones consecutivas de crecimiento mayores que 1.25, con corte duro
120. Son criterios numéricos finitos, distintos de probar convergencia o
explosión asintótica.

Se ejecutaron 17\,400 trayectorias: 7200 en radios \(r\leq0.01\) y
10\,200 en \(r=0.03,0.1,0.3\). No hubo contactos en las 7200 sondas locales
ni en las 6600 correspondientes a \(r=0.03,0.1\). En \(r=0.3\), entre
3600 sondas, se observaron 37 contactos: 19 desde el equilibrio positivo y
18 desde el negativo. Sus distancias iniciales reales al equilibrio
estuvieron entre \(0.1941805298\) y \(0.2997254735\). Por tanto, el primer
radio de bola muestreado con acceso fue 0.3; esta observación no localiza
una frontera continua de cuenca. Todos los conteos corresponden al mismo
convenio de reinicio.

\subsection{Chua con arctangente: control de Wu y candidato c590}
\label{bre:res:arctan}

La característica se fija como
\begin{equation}
 h(x)=a_1x+a_2\arctan(\rho x),\qquad
 \Caputo qx=\alpha[y-x-h(x)].
 \label{bre:res:arctan-convention}
\end{equation}
Así, el coeficiente lineal de \(x\) dentro del corchete es \(-(1+a_1)\).
En el control de Wu, \(a_1=0.4\); el número 1.4 corresponde a \(1+a_1\).
Los restantes parámetros son
\[
 \begin{aligned}
 q&=0.99,&(\alpha,\beta,\gamma)&=(8.4562,12.0732,0.0052),\\
 a_2&=-1.5585,&\rho&=1.
 \end{aligned}
\]
La recurrencia ADM local de orden cuatro, \(h=0.01\), \(T=100\) y descarte
de 50, desde \(X_{0,\pm}=\pm(13.8,0.7093,-19.8768)^{\T}\), produjo dos
respuestas simétricas acotadas, con frecuencia dominante \(0.5398920216\)
y entropías \((0.11607348,0.11106582,0.11095067)\)
\cite{Wu2023ArctanChua}. Este control reconstruye la recurrencia local;
la sección~\ref{bre:num:adm-local} demuestra por qué no se identifica con
la integración causal de memoria completa.

El candidato c590 utiliza la misma familia de campos, con
\begin{equation}
 \begin{aligned}
 q&=0.9999,&\alpha&=21.8493569066,\\
 \beta&=19.0818408409,&\gamma&=0.0073780120,\\
 a_1&=0.0422897934,&a_2&=-3.3367815123,\\
 \rho&=1.7984259333.
 \end{aligned}
 \label{bre:res:c590-parameters}
\end{equation}
La exploración inicial en orden entero evaluó 2400 casos globales y 1000
casos locales; el espectro variacional de la selección entera fue
\((0.4699043683,0.0010414594,-19.0491693716)\). Su semilla
\((7.6733768929,0.5079327671,-14.4903935267)^{\T}\) se utilizó para evaluar
por separado los órdenes \(0.9995,0.9998,0.9999,0.99995\), todos mediante
ABM de memoria completa, \(h=0.005\), \(T=300\) y descarte de 150.
No se transportó una historia entre órdenes distintos.

De la trayectoria con \(q=0.9999\) se extrajeron 16 estados entre
\(t\simeq150\) y 300. Cada estado definió un reinicio independiente. La
criba corta usó \(h=0.0025\), \(T=180\) y descarte de 90; el control largo
comparó \(h=0.0025,0.005,0.01\), con \(T=300\) y descarte de 150.
La selección, extraída cerca de \(t=240\), fue
\begin{equation}
 X_0^{\rm frac}=\begin{pmatrix}
 5.8642449791\\1.5847111486\\3.2155806478
 \end{pmatrix}.
 \label{bre:res:c590-seed}
\end{equation}
Permaneció acotada en los tres pasos; las medianas remuestreadas de la
prueba 0--1 fueron
\[
 (K_{0.0025},K_{0.005},K_{0.01})=(0.865843,0.885089,0.740058).
\]
Los dos pasos finos superaron el umbral de selección 0.8. El valor del paso
mayor muestra sensibilidad numérica y evita atribuir a ese control una
convergencia monótona no observada.

Las ecuaciones estacionarias de las componentes segunda y tercera dan
\(y=\gamma x/(\beta+\gamma)\) y
\(z=-\beta x/(\beta+\gamma)\). La primera se reduce a
\[
 \left(a_1+\frac{\beta}{\beta+\gamma}\right)x
       +a_2\arctan(\rho x)=0.
\]
Esta reducción también permite contar las raíces. Escriba
\(g(x)=\kappa_a x+a_2\arctan(\rho x)\), con
\(\kappa_a=a_1+\beta/(\beta+\gamma)>0\). La función es impar, y
\[
 \begin{aligned}
 g'(0)&=\kappa_a+a_2\rho<0,\\
 g''(x)&=\frac{-2a_2\rho^3x}{(1+\rho^2x^2)^2}>0\quad(x>0).
 \end{aligned}
\]
Además, \(g'(x)\to\kappa_a>0\) y \(g(x)\to+\infty\) cuando
\(x\to+\infty\). La derivada estrictamente creciente tiene un único cero
positivo; la función decrece desde \(g(0)=0\), alcanza un mínimo negativo
y luego crece hasta infinito. Existe exactamente una raíz positiva; por
imparidad, exactamente una negativa, además del origen. Sus valores producen
\[
 E_0=0,\qquad E_\pm=\pm\begin{pmatrix}
 4.6494055135\\0.0017970023\\-4.6476085112
 \end{pmatrix}.
\]
El origen es inestable. El margen angular de Matignon de los equilibrios
externos es \(7.42265\times10^{-4}>0\), por lo que son localmente estables
para el orden fijado. Su coexistencia con una trayectoria candidata exige
estudiar cuencas; la estabilidad local no decide el destino de
\eqref{bre:res:c590-seed}.

Sobre 30\,000 estados postransitorios, ninguna componente cumplió el filtro
de apariencia periódica. Las entropías espectrales fueron
\((0.394751,0.261922,0.315229)\); las fracciones de potencia dominante,
\((0.122163,0.290492,0.241937)\); y las derivas relativas de frecuencia,
\((0.941176,0.029412,0.029412)\). Las
Figs.~\ref{bre:res:fig-c590} y~\ref{bre:res:fig-c590-diagnostics} acompañan
estos indicadores. La conclusión es un candidato aperiódico, con indicios
numéricos de caos; los exponentes enteros no se transfieren automáticamente
al sistema de Caputo.

\begin{figure*}[t]
 \centering
 \includegraphics[width=0.96\textwidth]{chua_caputo_arctan.pdf}
 \caption{Trayectoria postransitoria de c590, \(q=0.9999\), \(h=0.005\),
 desde \eqref{bre:res:c590-seed}. El intervalo de observación es finito.}
 \label{bre:res:fig-c590}
\end{figure*}
\begin{figure*}[t]
 \centering
 \includegraphics[width=0.96\textwidth]{arctan_diagnosticos.pdf}
 \caption{Trayectoria postransitoria de c590 y reconstrucción de su
 primer armónico, con frecuencia \(f_1=0.67451\). La diferencia geométrica
 muestra la parte de la respuesta que no describe ese armónico. La
 clasificación dinámica utiliza además el control de paso y la prueba 0--1.}
 \label{bre:res:fig-c590-diagnostics}
\end{figure*}

El sondeo de cuenca utilizó superficies esféricas, con seis semiejes y
direcciones de Fibonacci, en lugar de bolas uniformes. Cada PVI usó
\(h=0.005\), \(T=180\) y descarte de 90. La referencia incluyó la nube y
su inversión central. La calibración dio
\(d_{90}^{\rm cal}=0.2950508245\) y umbral
\(\delta=0.5901016489\). La parada por equilibrio se activó desde
\(t=45\), con distancia \(<10^{-6}\), norma del campo \(<10^{-4}\) y
50 pasos consecutivos; el corte de divergencia fue \(\norm{X}_2=120\).
Las 8400 sondas con \(r\leq0.3\) dieron cero contactos: 8396 completaron el
horizonte y cuatro terminaron anticipadamente por el criterio de equilibrio.
En \(r=1\) hubo 22 contactos de 2400, todos desde \(E_0\); en \(r=2\),
588 de 2700, distribuidos como 289, 150 y 149 desde \(E_0,E_+,E_-\).
Los contactos macroscópicos describen acceso a la nube bajo el muestreo
superficial. La ausencia local de contactos no excluye todas las condiciones
de una vecindad abierta.

\subsection{MAVPD entero: ciclos, candidato caótico y puntos perpetuos}
\label{bre:res:mavpd}

El campo MAVPD se definió en \eqref{bre:geo:mavpd}. Se fijan
\(\delta=100\), \(\rho=200\), mientras \(\xi\) y \(\gamma\) seleccionan el
régimen \cite{matouk2023}. La parte lineal de Lur'e tiene
\(b=-\delta e_1\), \(c=e_1\) y \(\psi(u)=u^3\). La transferencia en la
convención \(W=c^{\T}(sI-P)^{-1}b\) es
\begin{equation}
 \begin{aligned}
 W(s)&=-\frac{\delta D(s)}{Q(s)},\qquad D(s)=s^2+\xi s+\rho,\\
 Q(s)&=(s-\delta\gamma)D(s)-\delta s.
 \end{aligned}
 \label{bre:res:mavpd-transfer}
\end{equation}
Como \(\cos^3\theta=(3\cos\theta+\cos3\theta)/4\), el primer armónico de
\((A\cos\theta)^3\) tiene amplitud \(3A^3/4\), y
\(N(A)=3A^2/4\). La condición de fase, después de eliminar el denominador
no nulo y poner \(z=\omega^2\), queda
\begin{equation}
 z^2+(\delta-2\rho+\xi^2)z+\rho(\rho-\delta)=0.
 \label{bre:res:mavpd-frequency}
\end{equation}
Para precisar la ganancia, el polinomio de \(P+kbc^{\T}\) se obtiene
reemplazando \(\gamma\) por \(\gamma-k\). Su parte real en
\(s=\ii\omega\) impone
\(\delta(k-\gamma)(\rho-z)=\xi z\); por tanto,
\[
 k=\gamma+\frac{\xi z}{\delta(\rho-z)},\qquad
 A=\sqrt{4k/3},
\]
para una raíz con \(z>0\), \(z\ne\rho\), \(k>0\). El modo normalizado es
\[
 v=\begin{pmatrix}
 1\\ k-\gamma+\ii\omega/\delta\\
 \rho/\delta-\ii\rho(k-\gamma)/\omega
 \end{pmatrix},\qquad Y_0(\theta)=A\operatorname{Re}(v e^{\ii\theta}).
\]
Así se fija la semilla sin escoger sus componentes por separado.

\subsubsection{Régimen periódico y comparación con PP}
Para \((\xi,\gamma)=(3.1,0.1)\), las ramas son
\[
 \begin{array}{c|rr}
 &\text{rama 0}&\text{rama 1}\\\hline
 \omega&10.5975230311&13.3447557342\\
 k&0.1397015948&0.3518790503\\
 A&0.4315886851&0.6849613618
 \end{array}
\]
y las semillas de fase cero son
\[
 Y_{0,0}=\begin{pmatrix}0.4315887\\0.0171348\\0.8631774\end{pmatrix},
 \qquad Y_{0,1}=\begin{pmatrix}0.6849614\\0.1725274\\1.3699227\end{pmatrix}.
\]
Se integraron ambas ramas con fases
\(0,\pi\) y 21 niveles de continuación. Cada nivel utilizó RK4 con
\(h=0.002\), transitorio de 20 y cola de 10. La primera rama alcanzó dos
ciclos internos simétricos; la segunda, el conjunto externo de la
Fig.~\ref{bre:res:fig-mavpd-periodic}. Un refinamiento DOP853 empleó
tolerancias \((2\times10^{-11},2\times10^{-13})\), paso máximo 0.005,
transitorio de 600 y ventana de 1200; registró 2433 retornos positivos a
\(y_2=0\).

En otro control de sección, DOP853 con eventos densos registró 405
cruces; el rango de los últimos retornos fue del orden de \(10^{-13}\).
El periodo fue \(0.493177986749\), con desviación estándar
\(1.25\times10^{-14}\). Un cálculo independiente dio
\(0.493178150753\), y la distancia normalizada entre nubes externas fue
\(0.00061096936\). El espectro finito
\((-0.0017614,-0.3887396,-51.1886634)\) y la mediana 0--1
\(K=0.000440085726\) respaldan un ciclo estable de periodo uno. La
clasificación cuasiperiódica de la fuente no se reproduce para estas
ecuaciones y estos controles.

El sondeo común, con \(h=0.005\), transitorio de 180 y cola de 80, empleó
108 condiciones: tres equilibrios, radios \(10^{-5},10^{-3},10^{-2}\) y
12 direcciones por radio. Hubo 54 destinos en cada ciclo interno y cero
contactos con el externo, coincidentes en ambos cálculos. Para
\(\xi=3.5\), conservando \(\gamma=0.1\), el mismo banco dio 24 contactos
con cada uno de los dos ciclos internos. Este control demuestra la
capacidad del sondeo para detectar acceso; no convierte el radio mínimo
ensayado en una prueba sobre radios arbitrariamente pequeños.

Los PP exactos de \eqref{bre:geo:mavpd-roots}, para \(\xi=3.1\), son
\[
 p_\pm=\pm\begin{pmatrix}
 0.1825741858350554\\0.01217161238900369\\0.1448421874291439
 \end{pmatrix}.
\]
Satisfacen \(J_f(p_\pm)f(p_\pm)=0\) y
\(\norm{f(p_\pm)}=3.44265186329548>0\). El experimento DOP853 desde esos
puntos, \(T=120\), tolerancias \(10^{-10}/10^{-12}\), alcanzó los ciclos
simétricos. En \(t\geq80\), los rangos de las componentes fueron
\((0.37943,0.05722,0.97545)\). El detector registró 74 y 75 máximos,
periodos medios \(0.53760,0.53730\) y coeficientes de variación
\(0.0361,0.0223\). Estos indicadores pertenecen al piloto desde PP y no al
periodo del ciclo externo. La Fig.~\ref{bre:res:fig-mavpd-pp} muestra que
la cancelación inicial de \(J_ff\) no se conserva a lo largo de la órbita.

\begin{figure*}[t]
 \centering
 \includegraphics[width=0.96\textwidth]{mavpd_periodico.png}
 \caption{MAVPD periódico, \(\xi=3.1\). Arriba: proyecciones del ciclo
 externo en los planos \(y_1y_2\) y \(y_1y_3\). Abajo: clasificación de las
 108 sondas, con 54 destinos en cada ciclo interno y ninguno en el externo,
 y continuación de la rama 1. Los paneles describen experimentos distintos
 dentro del mismo caso.}
 \label{bre:res:fig-mavpd-periodic}
\end{figure*}
\begin{figure*}[t]
 \centering
 \includegraphics[width=0.96\textwidth]{mavpd_pp.png}
 \caption{MAVPD entero desde los dos PP exactos. Las cruces indican las
 semillas; el diagnóstico de aceleración distingue su anulación inicial
 de la evolución posterior.}
 \label{bre:res:fig-mavpd-pp}
\end{figure*}

\subsubsection{Continuación hacia un candidato caótico}
Las dos ramas en \((\xi,\gamma)=(3.1,0.1)\) no superaron el cribado
finito \(\widehat\lambda_1>0.02\). Se transportó entonces la rama de menor
frecuencia desde \(\xi=3.10\) hasta 2.85 en pasos de \(-0.01\), y luego
se varió \(\gamma\) respecto de una frontera espectral calculada.
En \(E_\pm=(\pm\sqrt\gamma,0,\pm\sqrt\gamma)\), el polinomio del
Jacobiano es \(\lambda^3+b_1\lambda^2+b_2\lambda+b_3\), con
\[
 b_1=\xi+2\delta\gamma,\quad
 b_2=\rho-\delta+2\delta\gamma\xi,\quad
 b_3=2\delta\gamma\rho.
\]
Para una raíz \(\lambda=\ii\Omega\), \(\Omega>0\), sus partes real e
imaginaria imponen \(b_3-b_1\Omega^2=0\) y
\(\Omega(b_2-\Omega^2)=0\). Por tanto,
\(\Omega^2=b_2>0\) y \(b_1b_2=b_3\). Con \(g=2\delta\gamma\),
\begin{equation}
 \begin{split}
 b_1b_2-b_3
 &=(\xi+g)(\rho-\delta+\xi g)-\rho g\\
 &=\xi g^2+(\xi^2-\delta)g+\xi(\rho-\delta).
 \end{split}
 \label{bre:res:mavpd-hopf}
\end{equation}
En \(\xi=2.85\), la raíz alta da
\(\gamma_H=0.14380379839949112\). Es una frontera candidata de Hopf;
no se ha demostrado aquí transversalidad ni no degeneración de la
bifurcación. El cribado de desplazamientos seleccionó
\[
 \gamma=\gamma_H+0.010=0.15380379839949113.
\]
Se seleccionó el mayor desplazamiento ensayado con exponente válido
\(\widehat\lambda_1>0.5\), cero contactos preliminares, cero ambigüedades y
cero fallos. El valor \(\xi=2.85\) fue fijado como extremo del transporte,
no obtenido por una optimización conjunta.

La trayectoria estricta usó DOP853, \(T=900\), descarte de 300, muestreo
0.02, tolerancias \((2\times10^{-11},2\times10^{-13})\) y paso máximo 0.01.
El cálculo variacional separado usó descarte de 300, acumulación de 1200,
QR cada 0.5 y tolerancias \((2\times10^{-12},2\times10^{-14})\). Se obtuvo
\[
 \begin{aligned}
 \widehat\lambda_1&=0.7133503985,\quad
 \widehat\lambda_2=-0.0009058833,\\
 \widehat\lambda_3&=-20.9236249563,\quad D_{KY}=2.0340497651.
 \end{aligned}
\]
La estimación de Kaplan--Yorke usa
\(D_{KY}=2+(\widehat\lambda_1+\widehat\lambda_2)/|\widehat\lambda_3|\),
pues las dos primeras suman positivamente y la suma total es negativa.
Un control que perturbó el estado y modificó simultáneamente tolerancias,
paso máximo y ventanas dio \(\widehat\lambda_1=0.6830338121\).
La integración independiente produjo
\((0.6938035753,0.0017936287,-20.9005938240)\).
Los exponentes máximos positivos persisten en estos controles finitos.

La sección \(y_2=0\), con orientación positiva, retuvo 1174 cruces mediante
interpolación lineal entre muestras. Las medianas 0--1 de dos coordenadas
de retorno fueron \(0.9985162431\) y \(0.9984323397\). Sobre el flujo,
el resultado fue sensible al submuestreo: tres de diez pasos apoyaron la
clase caótica, uno la regular y seis fueron inconclusos. Por eso la
interpretación conjunta se apoya en retornos y exponentes, con su error
finito, y no en una sola prueba 0--1. La
Fig.~\ref{bre:res:fig-mavpd-chaos} reúne geometría y evolución variacional.
Un control EFORK con \(h=0.002\) alcanzó la misma nube calibrada.

Los 108 sondeos principales y cuatro dirigidos sobre ambos signos del
autovector inestable del origen, en radios \(10^{-7}\) y \(10^{-4}\),
sumaron 112 trayectorias. Todas terminaron en un equilibrio según el
clasificador; no hubo contactos con la nube candidata, ambigüedades ni
fallos numéricos. Los sondeos principales usaron radios
\(10^{-5},10^{-3},10^{-2}\), doce direcciones por equilibrio y radio,
transitorio de 400 y cola de 100. El resultado constituye evidencia
numérica de un candidato caótico compatible con ocultedad en ese banco;
la separación de vecindades continuas permanece sin certificar.

\begin{figure*}[t]
 \centering
 \includegraphics[width=0.96\textwidth]{mavpd_fases.pdf}\\[3pt]
 \includegraphics[width=0.76\textwidth]{mavpd_lyapunov.pdf}
 \caption{MAVPD, \(\xi=2.85\), \(\gamma=0.15380379839949113\).
 Arriba: proyecciones de la trayectoria postransitoria. Abajo: evolución de
 los exponentes variacionales finitos. Los controles de paso, integrador,
 retornos y condiciones iniciales delimitan la clasificación numérica.}
 \label{bre:res:fig-mavpd-chaos}
\end{figure*}

\subsection{Kalman--Fitts: obstrucción de ganancia y mapa de conmutación}
\label{bre:res:kalman}

El sistema de cuarto orden es
\begin{equation}
 \begin{gathered}
 \dot X=AX+b\tanh(c^{\T}X/\varepsilon),\\
 b=e_4,\quad c=-e_3,\quad\varepsilon=0.01,
 \end{gathered}
 \label{bre:res:kalman-field}
\end{equation}
con \(A\) compañera, última fila \((-a_0,-a_1,-a_2,-a_3)\), y polinomio
\[
 \begin{aligned}
 p_A(s)&=[(s+\beta)^2+m_1^2][(s+\beta)^2+m_2^2],\\
 (m_1,m_2,\beta)&=(0.9,1.1,0.03),\\
 (a_0,a_1,a_2,a_3)&=(0.98191881,0.121308,2.0254,0.12).
 \end{aligned}
\]
Las tres primeras filas de \(A\) desplazan consecutivamente las
coordenadas. La condición de equilibrio obliga a \(x_2=x_3=x_4=0\), y
la última ecuación da \(x_1=0\) porque \(a_0>0\)
\cite{kalman2019}. El Jacobiano en el origen tiene polinomio
\(s^4+a_3s^3+(a_2+1/\varepsilon)s^2+a_1s+a_0\). Sus coeficientes son
positivos, y las otras dos desigualdades de Routh--Hurwitz dan
\[
 \begin{aligned}
 a_3(a_2+1/\varepsilon)-a_1&=12.12174>0,\\
 a_3(a_2+1/\varepsilon)a_1-a_1^2-a_3^2a_0
 &=1.456324405056>0.
 \end{aligned}
\]
El único equilibrio es localmente asintóticamente estable.

Al resolver \((sI-A)v=b\), las tres primeras filas imponen
\(v_j=s^{j-1}v_1\), y la cuarta da \(p_A(s)v_1=1\). De aquí
\(W(s)=-s^2/p_A(s)\). La condición de fase da
\(\omega_0^2=a_1/a_3\), y la ganancia requerida es
\[
 \begin{aligned}
 \omega_0&=1.005435229142,\\
 k&=\frac{\omega_0^4-a_2\omega_0^2+a_0}{\omega_0^2}
   =-0.043168701157.
 \end{aligned}
\]
Para \(u=A\cos\theta/\varepsilon\), la función descriptiva es
\[
 N(A)=\frac1{\pi\varepsilon}\int_0^{2\pi}
         \frac{\tanh u}{u}\cos^2\theta\,\dd\theta.
\]
La extensión continua en \(u=0\) vale uno, y
\(0<\tanh u/u\leq1\); por ello \(N(A)>0\) para \(A>0\).
No existe amplitud que satisfaga la ganancia negativa. Esta obstrucción
justifica cambiar de localizador.

El precursor reemplaza \(\tanh\) por \(\operatorname{sign}\).
En un semiespacio con signo fijo \(\sigma\in\{-1,1\}\), la solución exacta
es
\begin{equation}
 \Phi_\sigma(t,X)=e^{At}X+A^{-1}(e^{At}-I)b\sigma.
 \label{bre:res:kalman-flow}
\end{equation}
Se obtiene integrando la fuerza constante:
\(\int_0^t e^{A(t-s)}\,\dd s=A^{-1}(e^{At}-I)\), con \(A\) invertible.
En \(\Sigma=\{c^{\T}X=0\}\), se elige el signo de salida y se busca el
primer cero transversal de \(c^{\T}\Phi_\sigma(t,X)\) después de abandonar
la sección. Si \(t_+(X)\) es ese tiempo, el mapa es
\(\Pi(X)=\Phi_\sigma(t_+(X),X)\). La transversalidad permite distinguir
el retorno de un contacto tangencial. El precursor convergió en 135
iteraciones a un retorno de periodo dos, \(\Pi^2(X)=X\), con semilla
\[
 X_{\rm sw}=\begin{pmatrix}
 0.212347571803\\-1.671123727247\\0\\1.478087383043
 \end{pmatrix}.
\]
Se continuó la no linealidad
\((1-\lambda)\operatorname{sign}\sigma+
\lambda\tanh(\sigma/\varepsilon)\) en 51 nodos
\(\lambda=0,0.02,\ldots,1\). Cada nodo utilizó RK4, \(h=0.01\),
transitorio de 100 y ventana retenida de 20. La corrida final descartó
300 unidades y conservó otras 300.

El ciclo resultante tuvo periodo \(11.9169960234\), frente a
\(11.9169379538\) en una integración independiente. Su espectro finito
fue
\[
 \widehat\lambda=\begin{pmatrix}
 -0.00283635\\-0.02423392\\-0.03056980\\-0.03988642
 \end{pmatrix},
\]
compatible con un exponente tangencial próximo a cero y tres negativos.
Los controles combinaron DOP853 con dos ajustes de tolerancias y RK4 con
\(h=0.01,0.005\), en horizontes 300 y 600. En las ocho configuraciones,
los errores relativos de periodo estuvieron entre \(4.87\times10^{-6}\)
y \(6.11\times10^{-6}\), y las distancias normalizadas entre nubes,
entre \(2.03\times10^{-4}\) y \(9.24\times10^{-4}\).

De las 48 condiciones alrededor del origen, doce en cada radio
\(10^{-5},10^{-4},10^{-3},10^{-2}\), cuarenta se clasificaron en el
equilibrio y ocho como acotadas no triviales; ninguna contactó el ciclo.
Ambos cálculos coincidieron punto por punto. La evidencia favorece un
ciclo estable compatible con ocultedad bajo ese sondeo. Una demostración
requiere además un dominio de retorno con atracción certificada y una
vecindad explícita del origen contenida en su cuenca local.

\subsection{PLL: retorno sobre el cilindro desde una solución exacta}
\label{bre:res:pll}

El PLL de \eqref{bre:geo:pll} tiene
\(\tau_1=0.0448\), \(\tau_2=0.0185\), \(L=500\) y
\(\omega_\Delta=178.9\) \cite{pll2015}. Para separar la constante del
filtro de los horizontes de integración, se escribe
\(T_f=\tau_1+\tau_2=0.0633\). Sus ecuaciones son
\begin{equation}
 \begin{aligned}
 \dot x&=-\frac{x}{T_f}+\frac{\tau_1}{2T_f}\sin\theta,\\
 \dot\theta&=\omega_\Delta-\frac L{T_f}x
                    -\frac{\tau_2L}{2T_f}\sin\theta.
 \end{aligned}
 \label{bre:res:pll-field}
\end{equation}
El estado pertenece a \(\R\times S^1\). Las condiciones estacionarias
imponen \(x_e=(\tau_1/2)\sin\theta_e\) y
\(\sin\theta_e=2\omega_\Delta/L=0.7156\). En una vuelta se obtienen
\[
 \begin{aligned}
 x_e&=0.01602944,\\
 \theta_f&=0.797482699881,\qquad
 \theta_s=2.344109953708.
 \end{aligned}
\]
El Jacobiano tiene
\[
 \operatorname{tr}J_e=-\frac{1+\tau_2L\cos\theta_e/2}{T_f},\qquad
 \det J_e=\frac{L\cos\theta_e}{2T_f}.
\]
El primer punto es un foco estable, con espectro
\(-33.41714171\pm40.52189661\ii\); el segundo es una silla, con espectro
\((-37.78074660,73.01945339)\).

Después de trasladar el equilibrio y conservar la no linealidad sinusoidal,
la transferencia es
\[
 \begin{aligned}
 W(s)&=-\frac L2\frac{1+\tau_2s}{s(1+T_fs)},\\
 \operatorname{Im}W(\ii\omega)
 &=\frac{L(1+\tau_2T_f\omega^2)}{2\omega(1+T_f^2\omega^2)}>0.
 \end{aligned}
\]
No hay cruce real para \(\omega>0\), de modo que el cierre directo con
una función descriptiva estática real no proporciona semilla.

Para \(L>0\), introducir \(y=\dot\theta\) permite despejar
\[
 x=\frac{T_f}{L}(\omega_\Delta-y)-\frac{\tau_2}{2}\sin\theta.
\]
Derivar \(y\), sustituir \(\dot x\) y agrupar
\(\tau_1+\tau_2=T_f\) da
\[
 T_f\dot y=\omega_\Delta-\frac L2\sin\theta
              -\left(1+\frac{\tau_2L}{2}\cos\theta\right)y.
\]
En una rotación positiva, \(y>0\) permite usar \(\theta\) como variable
independiente y dividir entre \(\dot\theta=y\):
\begin{equation}
 \frac{\dd y}{\dd\theta}=H_L(\theta,y)
 =\frac{\omega_\Delta-(L/2)\sin\theta}{T_fy}
    -\frac{1+(\tau_2L/2)\cos\theta}{T_f}.
 \label{bre:res:pll-angle}
\end{equation}
El mapa \(P_L(y_0)=y(2\pi;y_0)\) sólo se define en los datos que conservan
\(y>0\) durante toda la vuelta. Un punto fijo da una órbita corredora, con
periodo físico \(\mathcal T=\int_0^{2\pi}y(\theta)^{-1}\,\dd\theta\).
Si \(z=\partial y/\partial y_0\), la diferenciación del PVI produce
\[
 \begin{gathered}
 z'=-\frac{\omega_\Delta-(L/2)\sin\theta}{T_f y^2}z,\\
 z(0)=1,\qquad P_L'(y_0)=z(2\pi).
 \end{gathered}
\]
Por ello \(|P_L'|<1\) implica atracción transversal local del punto fijo
y \(|P_L'|>1\), repulsión. El criterio se aplica al retorno de la ODE
sobre el cilindro, no a un mapa de estados instantáneos de Caputo.

En \(L=0\), \eqref{bre:res:pll-angle} posee
\(y(\theta)\equiv\omega_\Delta\), con
\[
 \mathcal T_0=\frac{2\pi}{\omega_\Delta},\qquad
 P_0'(\omega_\Delta)=e^{-2\pi/(T_f\omega_\Delta)}.
\]
La reconstrucción que divide entre \(L\) no se usa en ese extremo.
En las variables originales, \(\theta(t)=\omega_\Delta t\), y proponer
\(x_p=A\sin\theta+B\cos\theta\) en la ecuación del filtro conduce a
\(A-\omega_\Delta T_fB=\tau_1/2\) y
\(\omega_\Delta T_fA+B=0\). Resolver estas dos ecuaciones da
\begin{equation}
 x_p(\theta)=\frac{\tau_1}{2[1+(\omega_\Delta T_f)^2]}
              (\sin\theta-\omega_\Delta T_f\cos\theta).
 \label{bre:res:pll-zero-gain}
\end{equation}
Las demás soluciones difieren en \(Ce^{-t/T_f}\). Su contracción durante
una vuelta coincide con \(P_0'\), lo que justifica utilizar esta órbita
exacta para iniciar la continuación en 101 valores hasta \(L=500\).
El retorno usó tolerancias \((2\times10^{-11},2\times10^{-13})\) y paso
angular máximo 0.02.

En la sección \(\theta=0\pmod{2\pi}\), se encontraron los resultados de
la Tabla~\ref{bre:res:tab-pll}. Los errores de retorno cilíndrico fueron
\(2.68\times10^{-15}\) para el ciclo estable y \(6.18\times10^{-14}\)
para la órbita inestable. La reproducción independiente dio periodos
\(0.0826766624969991\) y \(0.108541779641824\), con multiplicadores
\(0.817204143145181\) y \(1.44110733084579\).
Las diferencias angulares se calcularon módulo \(2\pi\).

\begin{table}[t]
 \centering
 \caption{Retornos del PLL en \(L=500\), con fase inicial cero.}
 \label{bre:res:tab-pll}
 \begin{tabular}{@{}lrr@{}}
 \toprule
 Magnitud & Ciclo estable & Órbita inestable\\\midrule
 \(y_0\) & 140.125084913 & 135.172047420\\
 \(x_0\) & 0.004908904250 & 0.005535958797\\
 \(\mathcal T\) & 0.0826766624974 & 0.1085417796382\\
 \(P_L'(y_0)\) & 0.8172041431535 & 1.4411073306967\\
 \bottomrule
 \end{tabular}
\end{table}
\begin{figure*}[t]
 \centering
 \includegraphics[width=0.90\textwidth]{pll_cilindro.pdf}
 \caption{PLL sobre el cilindro de fase. El retorno distingue una órbita
 corredora estable y una órbita periódica inestable que organiza la
 separación entre destinos. Las distancias respetan la periodicidad angular.}
 \label{bre:res:fig-pll}
\end{figure*}

Los sondeos utilizaron los dos equilibrios, cuatro radios
\(10^{-5},10^{-4},10^{-3},10^{-2}\) y doce direcciones por radio.
Las 96 condiciones terminaron en el foco y ninguna contactó el ciclo;
las clasificaciones coincidieron entre los dos cálculos. La órbita
inestable de la Fig.~\ref{bre:res:fig-pll} aporta una descripción de la
frontera que el balance armónico directo no ofrecía. Los valores numéricos
del retorno y los sondeos no sustituyen encierros intervalares ni una
prueba global de separación de cuencas.

\subsection{Sistemas fraccionarios no Chua}
\label{bre:res:nonchua-fo}

\subsubsection{Lorenz generalizado: estabilidad local y absorción}
El campo de \eqref{bre:geo:lorenz} corresponde a
\((\sigma,r,a)=(3.4,6.8,-0.5)\), con \(\sigma=-ar\), y se considera a
\(q=0.995\) \cite{Danca2017}. Para calcular los equilibrios, la tercera
ecuación da \(z=xy\). La segunda queda
\(rx-y-x^2y=0\), y como \(1+x^2>0\),
\[
 y=\frac{rx}{1+x^2},\qquad z=\frac{rx^2}{1+x^2}.
\]
Si \(x=0\), se obtiene el origen. Si \(x\ne0\), sustituir estas
expresiones en la primera ecuación y usar \(\sigma=-ar\) da, con
\(u=x^2>0\),
\[
 0=arx\left[1-\frac r{1+u}-\frac{ru}{(1+u)^2}\right].
\]
Como \(arx\ne0\), se puede dividir y multiplicar por \((1+u)^2\):
\begin{equation}
 \begin{aligned}
 0&=(1+u)^2-r(1+u)-ru\\
  &=u^2+2(1-r)u+(1-r),\\
 u_\pm&=r-1\pm\sqrt{r(r-1)}.
 \end{aligned}
 \label{bre:res:lorenz-equilibria-algebra}
\end{equation}
Para \(r>1\), \(\sqrt{r(r-1)}>r-1\), de modo que \(u_-<0\) se excluye
y \(u_+>0\) genera dos signos de \(x\). Son todos los equilibrios: el
origen y
\[
 E_\pm=\begin{pmatrix}
 \pm3.47564776512854\\\pm1.80689408468029\\6.28012738724303
 \end{pmatrix}.
\]
El Jacobiano es
\begin{equation}
 J_L=\begin{pmatrix}
 -\sigma&\sigma-az&-ay\\r-z&-1&-x\\y&x&-1
 \end{pmatrix}.
 \label{bre:res:lorenz-jacobian}
\end{equation}
En el origen, la tercera coordenada se desacopla, por lo que
\[
 p_0(s)=(s+1)[(s+\sigma)(s+1)-\sigma r].
\]
La sustitución de los equilibrios exactos en el campo lo anula; el residuo
máximo al evaluar sus coordenadas numéricas fue \(1.07\times10^{-14}\).
El origen tiene autovalores
\[
 \begin{aligned}
 \lambda_1&=-7.15580467734555,\\
 \lambda_2&=2.75580467734555,\qquad \lambda_3=-1,
 \end{aligned}
\]
y es inestable. En \(E_\pm\), los autovalores son
\[
 \begin{aligned}
 \lambda_1&=-5.37029951773135,\\
 \lambda_{2,3}&=-0.0148502411343235
                    \pm3.83248917224189\ii.
 \end{aligned}
\]
El margen angular a \(q=0.995\) es \(0.0117287914887680>0\), de modo que
ambos equilibrios externos son localmente estables.
La eliminación de la sección~\ref{bre:geo:section} proporciona además dos
semillas FPP--A, relacionadas por \((x,y,z)\mapsto(-x,-y,z)\).
En las variables \(u=x^2\), \(v=y/x\), una raíz tiene
\((u,v,z)\simeq(0.5564665577,-16.9736823118,-7.4423053062)\), y se
reconstruye mediante \(x=\pm\sqrt u\), \(y=vx\). Su norma de campo es
\(23.42210895339\). Su integración y clasificación de destinos siguen
pendientes; resolver \(J_ff=0\) no determina su cuenca.

Para este campo sí puede demostrarse una cota global, válida para cualquier
reinicio y para \(0<q\leq1\). Considérese
\begin{equation}
 V(x,y,z)=x^2+y^2+\tfrac12(z-20.4)^2.
 \label{bre:res:lorenz-V}
\end{equation}
Su Hessiana \(\operatorname{diag}(2,2,1)\) es definida positiva. En cada
intervalo compacto de existencia, \(X-X_0=I^qF(X)\), con \(F(X)\)
acotado. La desigualdad convexa \eqref{bre:fund:convex}, aplicada a
\(V(X_0+\xi)-V(X_0)\), justifica
\({}^{\mathrm C}D^q V(X)\leq\nabla V(X)\cdot F(X)\); no se utiliza una
regla de cadena fraccionaria con igualdad
\cite{Gomoyunov2018Convex}. El cálculo específico es
\begin{align*}
 \nabla V\cdot F
 & =2x[-3.4(x-y)+0.5yz]\\
 &\quad+2y[6.8x-y-xz]+(z-20.4)(-z+xy)\\
 &=-6.8x^2-2y^2-z^2+20.4z\\
 &=-V+208.08-5.8x^2-y^2-\tfrac12z^2\\
 &\leq -V+208.08.
\end{align*}
Los términos \(xyz\) se cancelan porque sus coeficientes suman
\(1-2+1=0\); los términos \(xy\), porque
\(6.8+13.6-20.4=0\). La constante es
\(20.4^2/2=208.08\).

Escribiendo \(v(t)=V(X(t))\) y
\(r(t)=208.08-v(t)-{}^{\mathrm C}D^qv(t)\geq0\), la fórmula lineal de
variación \eqref{bre:fund:variation} da
\begin{align*}
 v(t)&=v(0)E_q(-t^q)+208.08[1-E_q(-t^q)]\\
 &\quad-\int_0^t(t-s)^{q-1}E_{q,q}(-(t-s)^q)r(s)\,\dd s.
\end{align*}
El núcleo es no negativo. Para \(0<q<1\), esto se sigue de la
representación positiva \cite{Podlubny1999}
\[
 E_q(-t^q)=\frac{\sin(\pi q)}{\pi}
 \int_0^\infty
 \frac{e^{-ut}u^{q-1}}{u^{2q}+2u^q\cos(\pi q)+1}\,\dd u,
\]
pues al derivar para \(t>0\) se obtiene
\(-\frac{\dd}{\dd t}E_q(-t^q)=t^{q-1}E_{q,q}(-t^q)\geq0\).
Integrar esta identidad da
\(\int_0^t u^{q-1}E_{q,q}(-u^q)\,\dd u=1-E_q(-t^q)\), usado para
la fuerza constante. La representación implica también
\(0<E_q(-t^q)\leq1\) y \(E_q(-t^q)\to0\). Para \(q=1\), las mismas
propiedades corresponden a \(e^{-t}\). Por consiguiente,
\begin{equation}
 \begin{aligned}
 V(X(t))&\leq208.08+[V(X_0)-208.08]E_q(-t^q),\\
 V(X(t))&\leq\max\{V(X_0),208.08\}.
 \end{aligned}
 \label{bre:res:lorenz-bound}
\end{equation}
Los subniveles de \(V\) son compactos. La cota impide que una solución
maximal abandone todos los compactos en tiempo finito; la continuación de
Volterra con historia conservada establece existencia global y unicidad.
Para una familia acotada \(B\) de reinicios, sea
\(C_B=\sup_{X_0\in B}V(X_0)<\infty\). Entonces
\[
 \sup_{X_0\in B}V(X(t;X_0))
 \leq208.08+(C_B-208.08)_+E_q(-t^q).
\]
Dado \(\epsilon>0\), elegir un tiempo desde el cual el último término sea
menor que \(\epsilon\) prueba absorción uniforme por
\(\{V\leq208.08+\epsilon\}\). Esta conclusión controla todas las
trayectorias de reinicio; no identifica un atractor oculto ni demuestra
atracción en un espacio funcional admisible aún por construir.

\subsubsection{Rabinovich--Fabrikant: candidatos algebraicos}
El campo \eqref{bre:geo:rf} se evalúa con
\((a,b,q)=(0.1,0.2876,0.998)\) \cite{Danca2017}. La ecuación estacionaria
tercera es \(z(b+xy)=0\). Si \(z=0\), las primeras dos se escriben,
con \(u=x^2\), como
\[
 \begin{pmatrix}a&u-1\\1-u&a\end{pmatrix}
 \begin{pmatrix}x\\y\end{pmatrix}=0.
\]
Su determinante \(a^2+(u-1)^2>0\) exige \(x=y=0\); esta rama sólo da el
origen. Si \(z\ne0\), entonces \(xy=-b\ne0\), por lo que \(x\ne0\) y
\(y=-b/x\). La primera ecuación, multiplicada por \(x\), da
\(-b(z-1+u)+au=0\), de donde
\[
 z=1+\frac{a-b}{b}u.
\]
La segunda, también multiplicada por \(x\), es
\(u(3z+1-u)-ab=0\). Sustituir \(z\) y multiplicar por \(b\) produce
\begin{equation}
 \begin{aligned}
 &(3a-4b)u^2+4bu-ab^2=0,\\
 &u_\pm=\frac{b[-2\pm\sqrt{4+a(3a-4b)}]}{3a-4b}.
 \end{aligned}
 \label{bre:res:rf-equilibria-algebra}
\end{equation}
Para los parámetros elegidos, \(3a-4b=-0.8504\); el discriminante es
positivo y las raíces tienen suma y producto positivos. Ambas son
positivas y distintas. Cada una permite reconstruir dos equilibrios:
\(x=\pm\sqrt u\), \(y=-b/x\),
\(z=1+(a-b)u/b\). La sustitución verifica \(z\ne0\) y las tres ecuaciones.
Se obtienen exactamente cinco equilibrios:
\[
 \begin{gathered}
 E_0=(0,0,0),\\
 E_{1,2}=\begin{pmatrix}
 \mp1.15997695583801\\\pm0.247935959893466\\0.122306917444673
 \end{pmatrix},\\
 E_{3,4}=\begin{pmatrix}
 \mp0.085021329989515\\\pm3.38268055834300\\0.995284804098130
 \end{pmatrix}.
 \end{gathered}
\]
El Jacobiano es
\begin{equation}
 J_R=\begin{pmatrix}
 a+2xy&z-1+x^2&y\\
 3z+1-3x^2&a&3x\\
 -2yz&-2xz&-2(b+xy)
 \end{pmatrix}.
 \label{bre:res:rf-jacobian}
\end{equation}
En el origen, sus bloques dan
\(p_0(s)=(s+2b)[(s-a)^2+1]\) y autovalores
\(-2b,a\pm\ii\). El par complejo tiene orden crítico
\(2\arctan(1/a)/\pi=0.936548965138893\), menor que 0.998.
Así, el origen es inestable. La evaluación en \(E_{3,4}\) también da
inestabilidad, mientras \(E_{1,2}\) son estables, con margen angular
\(0.0435197213462810\). Los equilibrios exactos anulan el campo; el
residuo máximo de sus coordenadas numéricas fue \(1.06\times10^{-14}\). El conteo algebraico exhaustivo de
la sección~\ref{bre:geo:section} produce ocho FPP--A: cuatro en \(z=0\) y
cuatro fuera de ese plano. La realización multicanal permite formular el
balance de \eqref{bre:bal:mimo-balance}; ni sus ecuaciones ni las raíces
FPP--A reemplazan la continuación e integración del campo original.
Quedan pendientes los destinos de los cuatro pares de semillas, sus
controles de acotación y paso, y el sondeo de los cinco equilibrios.

\subsubsection{Muñoz--Pacheco: integración desde FPP--A}
El campo \eqref{bre:geo:munoz} tiene \(a=0.35\) y se evalúa a
\(q=0.97\) \cite{munoz2018}. No tiene equilibrios: \(f_2=0\) exigiría
\(|x|=1\), \(f_3=0\) impondría \(z=-xy\), y la primera componente
quedaría
\[
 f_1=-x(y^2-y+a)=-x\bigl[(y-1/2)^2+0.10\bigr]\ne0.
\] Las cuatro raíces suaves de
\eqref{bre:geo:munoz-roots} son
\[
 \begin{gathered}
 p_{1,2}=\begin{pmatrix}
 \pm0.816784043380077\\0.591607978309962\\\mp0.333568086760154
 \end{pmatrix},\\
 p_{3,4}=\begin{pmatrix}
 \pm3.183215956619923\\-0.591607978309962\\\mp5.066431913239846
 \end{pmatrix}.
 \end{gathered}
\]
Son condiciones iniciales algebraicas, separadas de la superficie no suave
\(x=0\). PECE con historia completa, \(h=0.01\) y \(T=35\), produjo las
cuatro trayectorias de la Fig.~\ref{bre:res:fig-munoz}. Las normas máximas
observadas fueron 6.421 para el par interior y 10.178 para el exterior.
Estos máximos sólo establecen finitud durante la ventana calculada.
La ausencia de equilibrios hace vacua la parte de la definición de ocultedad
que compara sus vecindades; siguen siendo necesarias existencia de un
conjunto invariante, atracción y acotación de la respuesta.

\begin{figure*}[t]
 \centering
 \includegraphics[width=0.96\textwidth]{munoz_semillas.png}
 \caption{Muñoz--Pacheco, \(q=0.97\), desde cuatro semillas FPP--A exactas,
 con \(h=0.01\) y \(T=35\). Las cruces marcan las condiciones iniciales.
 Las trayectorias son resultados preliminares de integración causal.}
 \label{bre:res:fig-munoz}
\end{figure*}

En el diagnóstico FPP--H de \eqref{bre:geo:history}, la comparación de
\(h=0.01\) y \(h=0.005\) sobre \([2,17.5]\) no produjo un mínimo
convergente a cero. El mínimo normalizado fino fue aproximadamente 0.0871
cerca de \(t=11.965\). El valor grueso 0.0254 en \(t=32.33\) queda fuera
del horizonte fino y no constituye una comparación de convergencia.
No se ha identificado un evento FPP--H. La familia direccional de frontera
\((0,y_0,0)\) tampoco ofrece un candidato acotado: sus soluciones exactas
son \(x=z=0\), \(y=y_0+t^q/\Gamma(q+1)\), como se dedujo en
\eqref{bre:geo:munoz-boundary}. La comparación de paso y horizonte de las
cuatro semillas suaves, y la reproducción del caso publicado a \(q=0.996\),
permanecen pendientes.

\subsubsection{Comparaciones de exponentes con otros parámetros}
Los ensayos anteriores no se confunden con dos comparaciones dinámicas de
referencia \cite{fischer2020}. En el Lorenz estándar
\[
 \begin{aligned}
 \Caputo qx&=\sigma(y-x),\\
 \Caputo qy&=\rho x-y-xz,\\
 \Caputo qz&=xy-\beta z,
 \end{aligned}
\]
con \((\sigma,\rho,\beta,q)=(10,200,8/3,0.985)\), el dato \(X_0=(0.1,0.1,0.1)\), \(h=0.001\), \(T=500\) e intervalo de
ortogonalización \(T_R=5\) dieron
\((-0.0027791532,-0.0868071712,-1.6224574503)\). La diferencia máxima
respecto de \((-0.0026,-0.0870,-1.6225)\) fue menor que
\(1.93\times10^{-4}\). La concordancia cuantitativa verifica ese cálculo
finito de referencia; no aporta un candidato caótico ni resultados sobre
el Lorenz generalizado anterior.

Para RF con \((a,b,q)=(0.1,0.98,0.999)\), desde el mismo dato inicial,
\(h=0.005\), \(T=300\) y \(T_R=0.02\), el espectro calculado fue
\((0.0611165030,0.0038981399,-1.8324646836)\), frente a
\((0.0749,0.0018,-2.0850)\). La diferencia del tercer exponente,
\(0.2525353\), excedió la tolerancia 0.05 y fue reproducida por otro
solucionador. El primer exponente positivo motiva estudiar la dinámica,
pero el ensayo no incluyó localización de una semilla ni sondeos de
vecindades. Por tanto, estos parámetros y sus diagnósticos no completan
la evaluación del candidato RF con \(b=0.2876\).
